\documentclass{article}
\usepackage{graphicx} % Required for inserting images
\usepackage{amsmath}
\usepackage[english, russian] {babel}
\usepackage[utf8]{inputenc}
\usepackage[T2A]{fontenc}
\usepackage{minted}
\usepackage{float}
\usepackage{amssymb}

\title{Языки программирования(Вольное слушиание)}
\author{silvia.lesnaia }


\begin{document}

\maketitle

\textbf{09.09.26}

\textbf{Пара №2}

Функциональное программрвоание предпологает, что вся проагрмма представляет из себя предсвталение функций. 

Функцию стоит редсталчть в математчиеском смысле. 

$f:\Delta \rightarrow z$

$f(x_1,.............,x_n) \rightarrow y$

Чистая функция - функция без побочных эффектов. В каких либо усолвиях йункция не взызываетя 
она дает чистый результат. Вызов функции не должен привоидиьь к имзенению окружению.

С точки зрения импертовррного прогармирвоание вызов фкнции не имеет смысла. 

Программа предлапоалегт что в ней только описание фунции, программе стоит задать вопрос и он заключтся в том что 
вызов фунции завит от определенных аргументов. 


Выражение - ключевое поняте для функлциольного прогарммироание. По сути 
функциальнлльне програмирование оперирует этим понятием.

Средсвта комбинирование - комбинение средств  - как сущесвтности определенное 
в языике програимоние можно состовет в более что то единое 

Срежвтсва асбтрации


$\lambda -$ исчлеслние Черча, где все описывается функциями


\section{Standard ML - академический язык программироания}

Функуия - вызов фунциия дает нам комбинирровать сущности. 

Выаржение состоятвтся из каких то данных достаочтно просто

Фкунцции задаются в ифникслной форме

$+ : real*real \rightarrow real$

$: int*int \rightarrow int$

$abs(25*a)$ и $(op+)(2,3) \Rightarrow S$ по сути разнице нет

REPL - позвоаляет обзаться с интопретатором


Можно использовать как спец ксотруки так и выражения

Список костнркуций: 
\begin{itemize}
    \item val a = 25*3 - дает обохнаечние - именеи обзоначем какоето выражение

     val a = abc - теперь новое знаечние ЕЕЕЕЕЕЕЕЕ ЗАТЕНЕНИЕ
    \item fun имя(параметры) : тип = выражение
\end{itemize}



Затенение - оно времнное или постоянное, оно может быть если исольузется обозначаение 
для другого обознчение

\textbf{10.09.26}

\textbf{Пара №3}


Хвостовая рекурсия - 

Лушче писать рекрсии то лучше писать их хвостовыми. 
В картеже обхясянются значения разного типа, т.е можно бъежнеить 
занечния и они могут принедалжеить к разным типам. 

(25,3,~48.3,"abc") - элементы картежа. 

Типы перечесиалются через * - int * int * real * string 

Пример: fun f(v : real * int) :  int * real * real = 

\#1(v) - функция вернет первый эелмент кортежа 

\#2(v) - фунция верет второй эелмент кортежа  

\# и число это название функция.  

(a,b) - это объеежнение картежа.

Записи - это как стуркутры в C++, т.е это наборым именовнонвых полей. 

Определние: val s =  ${name = "Sergei", surname = "Mironov", number = 345, summ = ~34.98}$

Имя поля и его знаенчие. Тип записи это тип сложный который явлется пересилсенем имен полей, 
с типами этих полей заклчюенный в фигурный скобки, причем порядко следование 
этих полей безразоличиен. 


\#s name вернет Sergei 
и тд 

list - список. Список бывает элементов определного типа. Нужно написать тип списка - int list 
-это список нормальных чисел. 

а (int * real) list это будет по сути чередвоание

int * real list полчуитсья список где одна часть номарьныз числа а другая вещесвтенная 

ОООО ДАААААА СИНТАКСИЧЕСКАИЙ САХАР

Это когда вовдится костркутция явелся упрощение записи другой более слоджной 
синтскичесйо конструкции.

hd l $\Rightarrow 2$ - выдаст голову списка 

hd [] - выдас ошибку ведь нельзя вывести голову пустого списка 

tl l $\Rightarrow$ [3,6,8,15] - хвост списка 

tl[5] $\Rightarrow$ [] - вернет пустой список ведь 5::[]

tl[] - выдаст ошибку 

null l $\rightarrow$ false - это фунция предикат, так как возварзет истинку или ложь

option - раюоатет только с чемто, как самостяотельного его нет. Это сталнартный тип 
данных в языках прогарммирование. Оно выдвает есть либо нет типа. 

У него есть выводы: NONE - ничего нет или SOME что то есть 

SOME 5 - типа int 

is SOME V - выдаст true если что то есть, т.е в опции есть занечние в противном случае 
false 

Фунция valOf(v) - если SOME что то вернется к примеру 5, если это NONE то выдаст ошибку. 
По сути просто мы распоковывеаем данные. 


\textbf{17.09.26}

\textbf{Пара №4}

Количество опрециий надо дорожиь. И использовать по возомжности дешевые операции. 

\subsection{Функция высших порядков}

Функции нисшего порядка не сущетсвует. Зато есть высшого лол. 

Это функия преинмает фкнцию в качестве аргумента. Т.е для функции аргментом 
выступает другая функция. 

Пример: val find = : int list * (int -> bool) -> int option

Функции врзрващающие выражениея это функции в прификсной форме. 


\end{document}